\subsection{Locating a citing passage} An opinion that left an issue open is
usually cited for something else, so a case-level edge is not evidence that a
later court said anything about the issue. For each origin we take its reporter
citations from the bulk citation table, build a pattern of the form
\emph{volume reporter page}, and search the text of each citing opinion for it.
Edges where no citation string can be located in the citing text are dropped
rather than guessed at; the count of such drops is reported. Around the located
citation we take a window of $\pm900$ characters, which is the citing passage.

\subsection{Deciding whether a passage is about the issue} A passage is treated
as issue-specific when at least three content words of the origin's issue
statement appear in it, after stopword removal. This is a filter on what gets
annotated, not a label. Annotators mark passages that are in fact about
something else \textsc{off\_issue}, and those are dropped, so false positives are
removed by annotation. False negatives are not recoverable: a later court that
restates the issue in different words will not be found. The rate of those is
unknown and is stated as a limitation.

\subsection{Capping per origin} At most three citing passages are drawn per
origin. Without a cap, a handful of heavily cited opinions would supply most of
the sample, and the clustered bootstrap would rest on a few very large clusters.

\subsection{The matched decided control} The comparison group is annotated
\textsc{decided} issues, drawn from the same benchmark and put through the same
chain construction. This is what separates status escalation from generic
citation compression. If later courts simply describe all prior holdings in
stronger and simpler terms than the originals warrant, both groups will show it.
The claim of interest requires escalation to be higher for issues the origin left
open, after matching on originating court, period, opinion length, and citation
volume.

\subsection{Inference} Several citing passages can share an origin, so passages
are not independent. All intervals for layer-two quantities are percentile
bootstrap resampling origins rather than passages, which is the level at which
the design clusters. We report descriptive contrasts with those intervals and
fit no causal model: whether a court leaves an issue open is a choice correlated
with how hard and how contested the issue is, and those same features plausibly
drive how later courts describe it.

\subsection{Directions through the citation network} Chains are not restricted to
the originating circuit. For each origin we record whether the citing opinion is
from the same court, another court of appeals, or the Supreme Court. These are
reported as descriptive pathways rather than as separate hypotheses.

\subsection{Blinding} The layer-two worksheet shows the origin's own passage and
never its label. An annotator told in advance that the origin left the issue open
would be judging the later court's language against a primed expectation, and the
central contrast in \S\ref{sec:long} is exactly between unresolved and decided
origins. Withholding the label costs nothing, since the origin passage carries
whatever information the judgement needs, and it removes the most obvious route
by which the contrast could be manufactured.
